\section{Introduction}

\label{sec:01}

\label{sec:introduction}

Electronic voting presents a persistent tension between public verification and ballot secrecy. Research on voting privacy distinguishes ballot secrecy from the various forms of verification that an election system can support \citep{jonker2013privacy}. Election participants need evidence that only eligible ballots are accepted, that votes are counted correctly, and that the announced result follows from the recorded election. Yet exposing each voter's selection can compromise the privacy on which a secret ballot depends. A blockchain can preserve a publicly inspectable sequence of election transactions, but the visibility of that record does not itself protect candidate choice. Pseudonymous identifiers conceal civil names only to the extent that they cannot be linked to their owners. Similarly, a contract that records valid transactions need not verify that the encrypted contents of those transactions represent admissible votes.

Biometric verification addresses an important but different requirement. Facial and fingerprint recognition can support identity-aware enrollment and help determine whether a person is authorized to vote. Fingerprint-based blockchain voting and hierarchical biometric-enabled election infrastructures have also been investigated \citep{ibrahim2021electionblock,indrason2025mbcsd}. These mechanisms do not necessarily prevent the authentication service or public ledger from observing ballot content. The distinction becomes particularly important when a biometric system is connected to a smart contract. Authentication may establish who can receive voting credentials, while a separate cryptographic process must define what the public verifier can learn about the resulting vote.

Our earlier study developed an IoT-assisted biometric voting architecture that combined fingerprint and facial verification, off-chain identity management, and blockchain vote processing \citep{rifat2026biometric}. That publication established the authentication and application foundation for the present research. Its voting pathway used candidate-identifying transaction data and maintained plaintext candidate totals. Those properties supported the earlier system's public recording model, but they did not provide encrypted ballot accumulation or threshold-based result disclosure. The previously reported biometric recognition and IoT timing results belong to that earlier evaluation. They are not new measurements in this study.

The privacy problem considered here arises after a voter has been authenticated. The system must permit the voter to cast a ballot without sending the selected candidate to the biometric authorization service or revealing it in the accepted voting transaction. At the same time, the contract must reject ballots that lack an eligible credential, reuse an election-scoped voting nullifier, select an invalid number of candidates, or contain ciphertexts inconsistent with the claimed ballot structure. Concealing plaintext selections without proving their validity would leave an important integrity requirement unresolved.

Existing cryptographic voting research offers essential foundations for this problem. Succinct zero-knowledge arguments can prove constraints on private data while exposing only a specified public statement \citep{groth2016snark}. Huber et al.\ investigated Groth16 proofs for the validity of exponential-ElGamal-encrypted ballots \citep{huber2025ballotvalidity}. Multi-authority encrypted elections and publicly checkable decryption also have an established research history \citep{cramer1997multiauthority}. More recent integrated architectures combine eligibility proofs, encrypted voting, and blockchain mechanisms \citep{alsuwat2026voting,jung2026anonymous}. Other systems incorporate non-interactive proofs, blind-signature mechanisms, and threshold encryption within blockchain elections \citep{wang2024versatile}. These contributions show that neither ballot-validity proofs nor homomorphic tallying are new primitives introduced by VoteChain-ZK.

The remaining question is how these mechanisms behave when connected to an existing biometric authorization application with identifiable voters, protected credentials, election administrators, and public blockchain observers. A proof of ballot validity does not establish that a biometric match was correct. Encryption does not conceal a candidate selection from a compromised client that sees the plaintext. A threshold decryption procedure may distribute key authority without removing the requirement to verify the published decryption evidence. The corresponding information boundaries must be specified alongside the cryptographic relations and the execution measurements.

This study presents VoteChain-ZK, an implemented extension that separates biometric authorization from private ballot preparation. The voter-facing web client initiates an authorization request that does not include candidate choice. Following a successful biometric check, the credential authorization service releases a protected voting credential and membership witness to the requesting client. A local privacy worker then prepares the candidate selection, constructs its ciphertexts, and generates the proof. The biometric service remains identity-aware, but the inspected authorization pathway does not require the selected candidate to cross that service interface.

The private voting path uses a separate Poseidon commitment tree and an election-specific smart contract. It must not be confused with the Semaphore-based registration and election anchor retained elsewhere in the application. The configured Groth16 circuit has Merkle depth 17 and 16 encrypted candidate slots. Its public statement binds the eligibility root, election context, election-scoped nullifier, and ciphertext coordinates. Its private witness contains the credential and membership path, one-hot candidate selection, and encryption randomness. The contract checks the proof and nullifier uniqueness before adding the ciphertexts to an encrypted tally. The accepted transaction does not contain a plaintext candidate identifier.

After voting ends, a three-out-of-five threshold procedure supports recovery of the aggregate candidate totals. Trustee partial decryptions and their consistency proofs are checked by an external transcript verifier. The contract records the finalized counts and a commitment to the verification transcript after applying its state and arithmetic conditions. The distinction matters because a successful finalization transaction does not by itself prove the correctness of the external decryption evidence. Independent auditability requires access to the transcript and verification of its relationship to the on-chain aggregate.

This investigation is organized around two research questions:

\noindent\textbf{RQ1---Candidate-blind authorization and encrypted-ballot validity:} How can identity-aware biometric authorization be separated from candidate selection while an election-specific on-chain proof enforces voter eligibility, nullifier uniqueness, candidate validity, and correct encrypted-ballot construction?

\noindent\textbf{RQ2---Verifiable tally recovery and measured execution:} How can encrypted ballots be accumulated and their final totals recovered through independently checkable threshold-decryption evidence, and what computational and transaction costs arise under the implemented fixed-circuit protocol and controlled workloads?

The contributions of this study are as follows:

\begin{enumerate}
\item \textbf{Candidate-blind biometric authorization and protected credential release.} We implement an authorization-to-ballot workflow that excludes candidate choice from the inspected biometric credential request and releases protected voting credentials and election-specific eligibility material after successful identity verification. Ballot construction takes place in a client-side privacy worker, with explicit recognition of credential-custody, client-integrity, and identity-to-nullifier linkage risks.

\item \textbf{Election-bound Groth16 validity and encrypted on-chain acceptance.} We implement a concrete voting relation combining private Poseidon eligibility membership, an election-scoped nullifier, frozen election-context binding, one-hot candidate constraints, and correctness of the 16 BabyJub exponential-ElGamal ciphertext slots. The Ethereum-compatible contract verifies its 73-signal public statement before accepting and homomorphically accumulating an encrypted ballot.

\item \textbf{Three-out-of-five threshold recovery and auditable finalization.} We implement encrypted tally accumulation and threshold recovery supported by externally validated distributed-key-generation and trustee discrete-logarithm equality evidence. The contract records final totals and a transcript commitment under state and arithmetic checks; independent review of the complete transcript, rather than the hash alone, is required to verify the decryption process.
\end{enumerate}

As supporting empirical evidence rather than an additional technical contribution, the performance study contains 2,000 measured accepted ballots across 25 independent elections. Their mean gas consumption was 1,314,307, and mean local submission-to-receipt latency was 760.640 ms. A separate fresh-chain campaign characterized the first accepted ballot. Candidate-count experiments used two, four, eight, and 16 active candidates with a fixed 16-slot representation. Eligible-set experiments covered 10 to 10,000 synthetic commitments while retaining Merkle depth 17. Mean Groth16 proving time ranged from 5.292 to 5.423 seconds under these configurations. These are local implementation measurements rather than public-network confirmation or end-to-end biometric voting times. Prior systematic research identifies scalability as an important concern in blockchain voting \citep{jafar2022scalability}; the present configurations do not establish network-level scaling behavior.

The resulting evidence supports the three stated technical contributions and an implemented means of verifying encrypted ballots without publishing their individual plaintext choices in the voting contract. It does not establish universal voter anonymity, receipt-freeness, coercion resistance, or production election readiness. The credential custodian can associate an identified voter with a public nullifier when it can recover the credential secret. Ballot confidentiality also depends on client integrity and trustee key protection. The experimental setup uses development-stage proof artifacts and a local Anvil network.

Section~\ref{sec:related-work} reviews the relevant biometric, cryptographic, and integrated voting literature. Section~\ref{sec:system-model} defines the system and trust assumptions. Sections~\ref{sec:architecture} and~\ref{sec:cryptographic-protocol} describe the architecture and cryptographic relation. Section~\ref{sec:implementation-experiments} presents the experimental design, followed by the results in Section~\ref{sec:results}. The discussion and limitations appear in Sections~\ref{sec:discussion-comparison} and~\ref{sec:limitations}, respectively. Section~\ref{sec:conclusion} concludes the study.